% QuickEIP Programming and Implementation Manual: shared preamble.
% Build: tectonic quickeip-manual.tex
\documentclass[11pt,twoside,openright]{book}

\usepackage[letterpaper,inner=1.15in,outer=1in,top=1in,bottom=1.05in,
            headheight=14pt,headsep=0.3in]{geometry}

% ---- fonts (XeTeX through Tectonic) ---------------------------------------
\usepackage{fontspec}
\usepackage{unicode-math}
\setmainfont{texgyrepagella}[
  Extension=.otf, UprightFont=*-regular, BoldFont=*-bold,
  ItalicFont=*-italic, BoldItalicFont=*-bolditalic, Ligatures=TeX]
\setsansfont{texgyreheros}[
  Extension=.otf, UprightFont=*-regular, BoldFont=*-bold,
  ItalicFont=*-italic, BoldItalicFont=*-bolditalic, Ligatures=TeX]
\setmonofont{DejaVuSansMono}[
  Extension=.ttf, UprightFont=*, BoldFont=*-Bold,
  ItalicFont=*-Oblique, BoldItalicFont=*-BoldOblique, Scale=0.84]
\setmathfont{texgyrepagella-math.otf}

\usepackage{microtype}
% Code words cannot hyphenate; let a line stretch instead of overflowing.
\setlength{\emergencystretch}{3em}

% ---- color, boxes, figures ------------------------------------------------
\usepackage{xcolor}
\definecolor{navy}{HTML}{1B2F4B}
\definecolor{steel}{HTML}{3D5A80}
\definecolor{teal}{HTML}{2A9D8F}
\definecolor{amber}{HTML}{C77700}
\definecolor{signal}{HTML}{B3261E}
\definecolor{noteblue}{HTML}{1F6FB2}
\definecolor{codebg}{HTML}{F5F7FA}
\definecolor{rulegray}{HTML}{C9D1DB}
\definecolor{commentgreen}{HTML}{4E7A55}
\definecolor{stringred}{HTML}{9A3324}

\usepackage{graphicx}
\usepackage{tikz}
\usetikzlibrary{positioning,arrows.meta,shapes.geometric,shapes.misc,calc,fit,backgrounds,decorations.pathreplacing}
\usepackage[most]{tcolorbox}

% ---- tables and lists -----------------------------------------------------
\usepackage{booktabs}
\usepackage{tabularx}
\usepackage{longtable}
\usepackage{array}
\usepackage{multicol}
\usepackage[shortlabels]{enumitem}
\newcolumntype{L}{>{\raggedright\arraybackslash}X}
\newcolumntype{P}[1]{>{\raggedright\arraybackslash}p{#1}}
\renewcommand{\arraystretch}{1.18}
\setlength{\tabcolsep}{5pt}
\setlist{itemsep=2pt,topsep=4pt}
\setlist[enumerate,1]{label=\arabic*.,leftmargin=*}

% ---- captions: Chicago style, label then period ---------------------------
\usepackage[font={small},labelfont={bf,sf},labelsep=period,
            justification=raggedright,singlelinecheck=false]{caption}
\captionsetup[table]{position=top,skip=4pt}
\captionsetup[figure]{position=bottom,skip=6pt}

% ---- headings -------------------------------------------------------------
\usepackage{titlesec}
\titleformat{\chapter}[display]
  {\sffamily\bfseries\color{navy}}
  {\normalsize\mdseries\scshape\color{steel}\chaptertitlename\ \thechapter}
  {4pt}
  {\Huge\raggedright}
  [\vspace{4pt}{\color{teal}\titlerule[1.2pt]}]
\titlespacing*{\chapter}{0pt}{-10pt}{26pt}
\titleformat{\section}{\sffamily\Large\bfseries\color{navy}}{\thesection}{0.7em}{}
\titleformat{\subsection}{\sffamily\large\bfseries\color{steel}}{\thesubsection}{0.7em}{}
\titleformat{\subsubsection}{\sffamily\normalsize\bfseries\color{steel}}{}{0em}{}
\titlespacing*{\section}{0pt}{16pt}{6pt}
\titlespacing*{\subsection}{0pt}{12pt}{4pt}
\titlespacing*{\subsubsection}{0pt}{10pt}{3pt}
\setcounter{secnumdepth}{2}
\setcounter{tocdepth}{1}

% ---- running heads --------------------------------------------------------
\usepackage{fancyhdr}
\pagestyle{fancy}
\fancyhf{}
\renewcommand{\headrulewidth}{0.4pt}
\renewcommand{\headrule}{{\color{rulegray}\hrule height\headrulewidth width\headwidth}}
\fancyhead[LE]{\sffamily\small\color{steel}\leftmark}
\fancyhead[RO]{\sffamily\small\color{steel}\rightmark}
\fancyfoot[LE,RO]{\sffamily\small\thepage}
\fancyfoot[RE,LO]{\sffamily\scriptsize\color{steel}QuickEIP Programming and Implementation Manual}
\renewcommand{\chaptermark}[1]{\markboth{\thechapter.\ #1}{}}
\renewcommand{\sectionmark}[1]{\markright{\thesection\ \ #1}}
\fancypagestyle{plain}{\fancyhf{}\renewcommand{\headrulewidth}{0pt}%
  \renewcommand{\headrule}{}\fancyfoot[LE,RO]{\sffamily\small\thepage}}

% ---- code listings --------------------------------------------------------
\usepackage{listings}
\lstdefinestyle{qe}{
  basicstyle=\ttfamily\small,
  keywordstyle=\color{navy}\bfseries,
  commentstyle=\color{commentgreen}\itshape,
  stringstyle=\color{stringred},
  backgroundcolor=\color{codebg},
  frame=leftline, framerule=2.5pt, rulecolor=\color{teal},
  xleftmargin=10pt, framexleftmargin=6pt, framesep=4pt,
  aboveskip=8pt, belowskip=8pt,
  columns=fullflexible, keepspaces=true, showstringspaces=false,
  breaklines=true, breakatwhitespace=true,
  upquote=true, tabsize=4,
}
\lstset{style=qe,language=C}
\lstdefinelanguage{shell}{morecomment=[l]{\#},sensitive=true}
\lstdefinelanguage{cmake}{
  morekeywords={cmake_minimum_required,project,list,include,add_subdirectory,
    add_executable,target_link_libraries,set,idf_component_register,
    idf_build_set_property,APPEND,PRIVATE,CACHE,BOOL,VERSION,SRCS,INCLUDE_DIRS,REQUIRES},
  morecomment=[l]{\#},morestring=[b]",sensitive=true}
\lstdefinelanguage{make}{morekeywords={include},morecomment=[l]{\#},sensitive=true}
\lstdefinelanguage{eds}{morecomment=[l]{\$},morestring=[b]",sensitive=true}
\lstdefinelanguage{yaml}{morecomment=[l]{\#},sensitive=true}

% ---- admonitions (ASD-STE100: warning, caution, note) ---------------------
\newtcolorbox{warning}{enhanced,breakable,colback=signal!5,colframe=signal,
  boxrule=0.8pt,arc=1.5pt,left=8pt,right=8pt,top=4pt,bottom=4pt,
  fonttitle=\sffamily\bfseries,coltitle=white,
  title={\raisebox{-1pt}{\large$\triangle$}\ \ WARNING}}
\newtcolorbox{caution}{enhanced,breakable,colback=amber!6,colframe=amber,
  boxrule=0.8pt,arc=1.5pt,left=8pt,right=8pt,top=4pt,bottom=4pt,
  fonttitle=\sffamily\bfseries,coltitle=white,title={CAUTION}}
\newtcolorbox{note}{enhanced,breakable,colback=noteblue!5,colframe=noteblue!80,
  boxrule=0pt,leftrule=3pt,arc=0pt,left=8pt,right=8pt,top=4pt,bottom=4pt,
  fontupper=\small}
\newcommand{\notelabel}{\textsf{\textbf{\color{noteblue}NOTE:}}\ }

% Procedure: a titled, numbered task (ASD-STE100 procedural writing).
\newtcolorbox{procedure}[1]{enhanced,breakable,colback=white,colframe=steel,
  boxrule=0.6pt,arc=1.5pt,left=6pt,right=8pt,top=4pt,bottom=4pt,
  fonttitle=\sffamily\bfseries\small,colbacktitle=steel!12,coltitle=navy,
  title={Procedure: #1}}

% API entry: name, then the prototype in a shaded box.
\newtcolorbox{protobox}{enhanced,colback=codebg,colframe=rulegray,boxrule=0.4pt,
  arc=1pt,left=6pt,right=6pt,top=3pt,bottom=3pt,before skip=3pt,after skip=4pt,
  fontupper=\ttfamily\small}
% Label keys must not contain an active underscore: the aux file of an
% \include'd chapter is read back after the underscore package activates it.
\ExplSyntaxOn
\str_new:N \l_qe_key_str
\cs_new_protected:Npn \qe_setkey:n #1
  { \str_set:Nn \l_qe_key_str {#1} \str_replace_all:Nnn \l_qe_key_str { _ } { - } }
\NewDocumentCommand{\apilabel}{m}{ \qe_setkey:n {#1} \exp_args:Nx \label { api: \l_qe_key_str } }
\NewDocumentCommand{\apipageref}{m}{ \qe_setkey:n {#1} \exp_args:Nx \pageref { api: \l_qe_key_str } }
\ExplSyntaxOff
\newcommand{\api}[2]{%
  \par\addvspace{10pt}\needspace{6\baselineskip}%
  \phantomsection\apilabel{#1}%
  \noindent{\sffamily\bfseries\color{navy}#1}\par\nopagebreak
  \begin{protobox}\raggedright #2\end{protobox}\nopagebreak}
\newcommand{\apiidx}[1]{\noindent\texttt{#1}\dotfill\apipageref{#1}\par}
\usepackage{needspace}

% ---- inline helpers -------------------------------------------------------
\newcommand{\code}[1]{\texttt{#1}}
\newcommand{\file}[1]{\texttt{#1}}
% Direction arrows that work in text and in math (subscripts of equations).
\DeclareRobustCommand{\OT}{\ifmmode\text{O}{\to}\text{T}\else O\kern0.05em$\to$\kern0.05em T\fi}
\DeclareRobustCommand{\TO}{\ifmmode\text{T}{\to}\text{O}\else T\kern0.05em$\to$\kern0.05em O\fi}
\newcommand{\hex}[1]{\texttt{0x#1}}
\newcommand{\term}[1]{\textit{#1}}
\newcommand{\menu}[1]{\textsf{#1}}

\usepackage{amsmath}
\usepackage{siunitx}
\sisetup{group-separator={,},group-minimum-digits=4}

% ---- links last -----------------------------------------------------------
\usepackage[hidelinks,bookmarksnumbered,pdfusetitle]{hyperref}
\hypersetup{pdftitle={QuickEIP Programming and Implementation Manual},
  pdfauthor={Adam G. Sweeney},pdfsubject={EtherNet/IP scanner and adapter stack},
  pdfkeywords={EtherNet/IP, CIP, embedded, C99, QuickEIP}}
\usepackage{bookmark}

% Load last: an active underscore breaks PGF if it loads first.
\usepackage[strings,nohyphen]{underscore}
